\scanpage{35}
b) filtration par ``\deleted{la}\unsure{le|la} \deleted{dimension}'' ``type dimensionnel''. % L1

On pose pour $i\geq 0$ % L2
\[\mathcal D_i(\mathcal M^+(X))=\deleted{\unsure{\mathcal M|\mathcal M^i}}\,\text{sous-catégorie pleine}\] % L3
de $\mathcal M^+(X)$ formée des objets % L4
qui se \inserted{loc. étale} dévissent [(du moins % L5
fibre par fibre)] en objets % L6
\margin{\deleted{\unsure{$M$|$\mathrm{Ob}$}}} % L6a
de la forme % L7
\deleted{$\unsure{M}$} $M(-j)$, avec $M\in\mathrm{Ob}\,\mathcal M^{+i}(X)$, % L8
\inserted{$j\in\mathbb Z,\ j\geq 0$} % L8a
i.e. on prend (du moins pour $X=\mathrm{Spec}\ k$ d'un % L9
corps) la filtration minima qui % L10
majore celle par le poids, \deleted{\unsure{et}} qui soit % L11
stable par $\otimes\deleted{\mathbb Z}\mathbb Q(-1)$, et \uline{épaisse}. % L12

Les $\mathcal D_i(\mathcal M^+(X))$ sont stables par l'image inverse, % L13
et pour l'image directe on a % L14
\[\begin{aligned}\mathbb R^j f_!(\mathcal D_i(\mathcal M^+(X)))&\subset\mathcal D_{i+j}(\mathcal M^+(Y))\\ % L15
\mathbb R^j f_*(\mathcal D_i(\mathcal M^+(X)))&\subset\mathcal D_{i+j}(\mathcal M^+(Y))\ \text{\textcircled{?}}\end{aligned}\] % L16
% margin note written sideways
\margin{\deleted{\unsure{le cas critique|les cas critiques} \unsure{est celui d'une} immersion ouverte}} % L16a
% slanted
\margin{cela distingue formellement les filtrations $\mathcal D_i(M)$ et $M_i$ !} % L16b
\inserted{[On a seulement $\mathbb R^jf_*(\mathcal M^{+(i)}(X))\subset\mathcal M^{+(i+2j)}(Y)$ !?]} % L17
% slanted
\margin{\unsure{tout reformuler|tout se formule|toute formule} en termes de la dim. de $X$ !} % L18

Ici \deleted{\unsure{encore}} on a % L19
\[M\in\mathrm{Ob}\,\mathcal D_i(\mathcal M^+(X))\Longleftrightarrow M(-j)\in\deleted{\unsure{\mathrm{Ob}}}\mathcal D_{i\deleted{\unsure{+j}}}(\mathcal M^+(X))\] % L20
\margin{pas $\llap{\uuline{\phantom{pas}}\kern.25em}\mathcal D_{i+j}$ !!} % L20a
et la filtration \inserted{de $\mathcal M^+(X)$} par les $\mathcal D_i(\mathcal M^+(X))$ s'étend % L21
en une filtration de $\mathcal M(X)$ par des $\mathcal D^i(\mathcal M(X))$. % L22

% archive sheet number: 34
